
This thesis presents the design and implementation of BlockQuiz, a web-based learning platform
for block-based programming exercises aimed at children aged 8--12. The work is motivated by the increasing importance of
computational thinking in school curricula and, in particular, by the Austrian subject "Digitale Grundbildung",
which creates demand for classroom-ready tools
that offer guided learning tasks, immediate feedback, and a deployment model suitable for schools.

Existing systems such as Scratch, Code.org, Microsoft MakeCode, and Blockly Games demonstrate the educational value of
block-based programming, but
they do not combine all of the requirements that are central to this project.
Some of them focus primarily on open-ended creativity, others on hardware programming or cloud-hosted delivery, and only a
subset support tightly constrained
exercises that can be assessed automatically.
BlockQuiz addresses this gap through a teacher-oriented platform for authoring and delivering short programming tasks with
clear learning goals, progressive hints, and
deterministic auto-grading.

The implementation is based on SvelteKit and Google Blockly. It includes a client-side sandbox for
running generated code, an authoritative server-side path that re-grades every submission inside an
isolated subprocess, a typed exercise model with publish validation, a grading framework for
input/output, turtle, and robot exercises, multilingual content, a teacher content-management system,
and identity-provider-synchronised classes with optional OIDC and SAML Single Sign-On. The whole
system follows a privacy-first deployment model that supports self-hosting in school networks.

The system was evaluated along two axes. A technical evaluation with an automated test suite and a
database verification script confirms that the exercise engines, the grading layer and the submission
pipeline behave correctly and that the implementation meets its requirements. A user evaluation with
two pre-service teachers found the platform understandable and pedagogically appropriate for its
target group and surfaced concrete usability improvements, several of which were addressed in
response. The results indicate that a focused, privacy-conscious block-programming platform with
automatic grading is both technically viable and positively received by the teachers expected to
author its content.
